% =============================================================================
% Appendix B — Expansion RAM (XRAM)
% =============================================================================
\chapter{Expansion RAM (XRAM)}

\epigraph{``More memory is always the answer.  The question is
irrelevant.''}{\textit{--- Every programmer, eventually}}

\section*{Introduction}

The NovaVM provides 512\,KB of expansion RAM (XRAM) --- eight times the 6502's
native 64\,KB address space.  XRAM is not directly addressable by the CPU;
instead, it is accessed through the Expansion Memory Controller (XMC) registers
at \$BA00--\$BA3F, four memory-mapped windows at \$BC00--\$BFFF, or via DMA
transfers (space~5).

This appendix provides a consolidated reference for XRAM's organization, access
methods, and named-block allocation system.  For the XMC register-by-register
documentation, see Chapter~10.


% ═══════════════════════════════════════════════════════════════════════════════
\section{Bank Structure}
\label{sec:xram-banks}
% ═══════════════════════════════════════════════════════════════════════════════

XRAM uses 24-bit addressing: an 8-bit bank number (high byte) combined with a
16-bit offset within the bank.

\begin{center}
\small
\begin{tabular}{@{}clll@{}}
\toprule
\textbf{Bank} & \textbf{Start Address} & \textbf{End Address} & \textbf{Size} \\
\midrule
0 & \$000000 & \$00FFFF & 64\,KB \\
1 & \$010000 & \$01FFFF & 64\,KB \\
2 & \$020000 & \$02FFFF & 64\,KB \\
3 & \$030000 & \$03FFFF & 64\,KB \\
4 & \$040000 & \$04FFFF & 64\,KB \\
5 & \$050000 & \$05FFFF & 64\,KB \\
6 & \$060000 & \$06FFFF & 64\,KB \\
7 & \$070000 & \$07FFFF & 64\,KB \\
\midrule
  & \multicolumn{2}{l}{\textbf{Total}} & \textbf{512\,KB} \\
\bottomrule
\end{tabular}
\end{center}

The active bank is selected by writing to \texttt{XmcBank} at \addr{BA0C}.
The total number of available banks can be read from \texttt{XmcBanks} at
\addr{BA0D}.


% ═══════════════════════════════════════════════════════════════════════════════
\section{Access Methods}
\label{sec:xram-access}
% ═══════════════════════════════════════════════════════════════════════════════

There are four ways to read and write XRAM, each suited to different use cases.

\subsection{Byte Commands (GetByte / PutByte)}

The simplest access method: set the bank, address, and data registers, then
issue a GetByte or PutByte command.  Good for reading or writing individual
values, but too slow for large transfers.

\begin{lstlisting}
10 REM write byte $42 to XRAM bank 0, offset $1000
20 XBANK 0
30 XPOKE $1000, $42
40 PRINT "READ BACK:"; XPEEK($1000)
\end{lstlisting}

\subsection{Bulk Commands (Stash / Fetch)}

Copy a block of bytes between CPU RAM and XRAM in a single operation.  The XMC
handles the transfer internally at full speed.

\begin{lstlisting}
10 REM copy 256 bytes from CPU RAM at $2000
20 REM into XRAM bank 0 at offset $0000
30 STASH $2000, 256, 0, 0
40 REM
50 REM copy them back to CPU RAM at $3000
60 FETCH $3000, 256, 0, 0
\end{lstlisting}

\subsection{XMC Windows (\$BC00--\$BFFF)}

Four 256-byte windows are mapped directly into the 6502 address space.  Each
window can be independently pointed at any 256-byte page in XRAM using the
XMAP command.  Once mapped, ordinary PEEK/POKE operations read and write XRAM
through the window.

\begin{center}
\small
\begin{tabular}{@{}cll@{}}
\toprule
\textbf{Window} & \textbf{CPU Address Range} & \textbf{Size} \\
\midrule
0 & \$BC00--\$BCFF & 256 bytes \\
1 & \$BD00--\$BDFF & 256 bytes \\
2 & \$BE00--\$BEFF & 256 bytes \\
3 & \$BF00--\$BFFF & 256 bytes \\
\bottomrule
\end{tabular}
\end{center}

\begin{lstlisting}
10 REM map window 0 to XRAM bank 0, page 16 ($1000)
20 XMAP 0, 0, 16
30 REM now PEEK/POKE at $BC00-$BCFF accesses XRAM
40 POKE $BC00, $42
50 PRINT PEEK($BC00)
60 XUNMAP 0 : REM release window
\end{lstlisting}

\begin{tipbox}
XMC windows are ideal for random access to small data structures stored in XRAM,
such as lookup tables or tile maps.  For sequential bulk transfers, Stash/Fetch
or DMA are faster.
\end{tipbox}

\subsection{DMA Transfers (Space 5)}

The DMA controller treats XRAM as space~5.  It can copy data between XRAM and
any other DMA space (CPU RAM, VGC character/color/graphics/sprite memory, or
even another region of XRAM).  This is the fastest method for large transfers.

\begin{lstlisting}
10 REM DMA copy 1000 bytes from XRAM offset $0000
20 REM to CPU RAM at $2000
30 DMACOPY 0, $2000, 1000, 5, 0
\end{lstlisting}


% ═══════════════════════════════════════════════════════════════════════════════
\section{Named Block Allocation}
\label{sec:xram-named}
% ═══════════════════════════════════════════════════════════════════════════════

The XMC provides a simple named-block filesystem within XRAM.  Each block has a
name (up to 27 characters) and can be stored and retrieved without needing to
track offsets manually.

\begin{center}
\small
\begin{tabular}{@{}lp{24em}@{}}
\toprule
\textbf{Command} & \textbf{Description} \\
\midrule
NStash  & Store a block of CPU RAM into XRAM under a given name. \\
NFetch  & Retrieve a named block from XRAM into CPU RAM. \\
NDelete & Delete a named block and free its XRAM pages. \\
NDirOpen / NDirRead & Enumerate all named blocks. \\
\bottomrule
\end{tabular}
\end{center}

\begin{lstlisting}
10 REM store level data under the name "LEVEL1"
20 XMEM "LEVEL1", $2000, 4096, 1
30 REM
40 REM retrieve it later
50 XMEM "LEVEL1", $2000, 0, 0
60 REM
70 REM list all named blocks
80 XDIR
90 REM
100 REM delete a named block
110 XDEL "LEVEL1"
\end{lstlisting}

\begin{notebox}
Named blocks are allocated from the XRAM page pool.  The NStash command
automatically allocates the required number of 256-byte pages.  Use XFREE to
check how many pages remain available.
\end{notebox}


% ═══════════════════════════════════════════════════════════════════════════════
\section{Page Usage Tracking}
\label{sec:xram-pages}
% ═══════════════════════════════════════════════════════════════════════════════

XRAM is divided into 256-byte pages.  With 512\,KB of XRAM, there are 2,048
pages in total.  The XMC tracks which pages are in use and which are free.

\begin{center}
\small
\begin{tabular}{@{}lll@{}}
\toprule
\textbf{Register} & \textbf{Address} & \textbf{Description} \\
\midrule
XmcPagesUsed & \addr{BA0E} (16-bit) & Number of pages currently allocated \\
XmcPagesFree & \addr{BA10} (16-bit) & Number of pages available \\
\bottomrule
\end{tabular}
\end{center}

The XALLOC command reserves a contiguous block of pages and returns the starting
offset.  XRESET releases all allocations and clears XRAM.

\begin{lstlisting}
10 PRINT "PAGES FREE:"; XFREE
20 XALLOC 16 : REM allocate 16 pages (4KB)
30 PRINT "PAGES FREE:"; XFREE
40 XRESET : REM free everything
50 PRINT "PAGES FREE:"; XFREE
\end{lstlisting}

\begin{warningbox}
XRESET clears all XRAM contents and allocations, including named blocks.  There
is no undo.  Use it only when you want to start fresh.
\end{warningbox}
